\subsection{The Case of Abelian Groups}

In this subsection, we assume that the group \(\mathrm{G}\) is abelian.

By Schur's lemma (VIII, p.~48, Corollary 1), every simple representation of G has dimension 1. Let \((\mathrm{M},\pi)\) be such a representation and \(\chi\) be its character; for every \(g\in G\) and \(x\in M\) , we have \(\pi(g)(x)=\chi(g)x\) . Consequently, the character \(\chi\) is a homomorphism from G to the multiplicative group \(K^{*}\) of K. Conversely, every homomorphism from G to \(K^{*}\) is the character of a representation of G of degree 1. So the set \(\widehat{G}\) of classes of simple K{[}G{]}-modules can be identified with the set \(\operatorname{Hom}(G,K^{*})\) of homomorphisms from G to \(K^{*}\) . We deduce from this an abelian group structure on \(\widehat{G}\) ; the product in \(\widehat{G}\) corresponds to the tensor product of representations. The groups G and \(\widehat{G}\) have the same cardinal by Proposition 5 of VIII, p.~411. Every function on G is central, and \(\widehat{G}\) is a basis of the vector space of mappings from G to K (loc. cit.). Because of the orthogonality relation for characters, such a mapping f has the following unique decomposition with respect to the basis \(\widehat{G}\) :

\[
f = \sum_ {\chi \in \widehat {\mathrm{G}}} \langle \chi , f \rangle_ {\mathrm{G}} \chi .\tag{39}
\]

For mappings \(f\) and \(f'\) from \(G\) to \(K\) , we have the relation

\[
\langle f, f ^ {\prime} \rangle_ {\mathrm{G}} = \sum_ {\chi \in \widehat {\mathrm{G}}} \langle \chi , f \rangle_ {\mathrm{G}} \langle \chi , f ^ {\prime} \rangle_ {\mathrm{G}}.\tag{40}
\]

Let \((\mathrm{V},\pi)\) be a linear representation of G. For any \(\chi\in\widehat{\mathrm{G}}\) , denote by \(V^{\chi}\) the subspace of V consisting of the vectors v such that \(\pi(g)(v)=\chi(g)v\) for every \(g\in\mathrm{G}\) . The space \(V^{\chi}\) is the isotypical component of type \(\chi\) of the K{[}G{]}-module V. The space V is the direct sum of the family \((V^{\chi})_{\chi\in\widehat{\mathrm{G}}}\) , and the projector \(p_{\chi}\) of V with image \(V^{\chi}\) associated with this decomposition is given by

\[
p _ {\chi} = | \mathrm{G} | ^ {- 1} \sum_ {g \in \mathrm{G}} \chi (g ^ {- 1}) \pi (g)\tag{41}
\]

because of relation (19).

Remark. --- Let n be the cardinal of the group G, and let \(\mu_{n}(\mathrm{K})\) be the group of n-th roots of unity in K. For every \(g \in G\) , we have \(g^{n} = 1\) ; consequently, \(\widehat{G}\) can be identified with the group \(\operatorname{Hom}(\mathbf{G}, \mu_{n}(\mathbf{K}))\) . The group \(\mu_{n}(\mathbf{K})\) is cyclic of order n (V, §11, No.~2, p.~78, Theorem 1). The group \(\widehat{G}\) is therefore isomorphic to the group \(\mathrm{D}(\mathrm{G}) = \operatorname{Hom}(\mathrm{G}, \mathbf{Q}/\mathbf{Z})\) . By VII, §4, No.~9, p.~26,

Proposition 10, the group \(\widehat{\mathbf{G}}\) is isomorphic to the group \(\mathbf{G}\) , and the mapping that sends an element \(g\) of \(\mathbf{G}\) to the homomorphism \(\chi \mapsto \chi(g)\) from \(\widehat{\mathbf{G}}\) to \(\mathrm{K}^*\) is an isomorphism from \(\mathbf{G}\) to \(\widehat{\widehat{\mathbf{G}}}\) .
% semantic-fragments: legacy-input-seam
\relax{}
